\chapter*{记号索引}
\phantomsection
\addcontentsline{toc}{chapter}{记号索引}

引用编号依次表示章、节、小节（在个别情形为习题）。

第 III 章 :

\(\mathcal{C}(\mathrm{X};\mathrm{E}),\mathcal{C}(\mathrm{X}),\mathcal{C}(\mathrm{X},\mathrm{A};\mathrm{E}),\mathcal{K}(\mathrm{X};\mathrm{E}),\mathcal{K}(\mathrm{X}),\mathcal{K}(\mathrm{X},\mathrm{A};\mathrm{E}),\mathcal{K}(\mathrm{X},\mathrm{A}),\mathcal{K}_{+}(\mathrm{X})\) （X 为局部紧空间，E 为拓扑向量空间）：III, 1, 1. Supp(f)（f 为取值于向量空间或 \(\overline{\mathbf{R}}\) 中的函数）：III, 1, 1。 \(\mathcal{C}^b (\mathrm{X};\mathrm{E}),\mathcal{C}^0 (\mathrm{X};\mathrm{E}):III,1,2.\) \(\| \mathbf{f}\|\) （f 为取值于赋范空间的函数）：III, 1, 2。 \(\mu (f),\langle f,\mu \rangle ,\int fd\mu ,\int f\mu ,\int f(x)d\mu (x),\int f(x)\mu (x)\) （f 为 \(\mathcal{H}(\mathrm{X};\mathbf{C})\) 中的函数， \(\mu\) 为（复）测度）：III, 1, 3。 \(\mathcal{M}(\mathrm{X};\mathbf{C}),\mathcal{M}(\mathrm{X}),\mathcal{M}_{\mathfrak{S}}(\mathrm{X};\mathbf{C}),\mathcal{M}_{\mathfrak{S}}(\mathrm{X}):III,1,3.\) \(\varepsilon_{a}:III,1,3.\) \(g\cdot \mu\) （g 为 \(\mathcal{C}(\mathrm{X};\mathbf{C})\) 中的函数）：III, 1, 4。 \(\overline{\mu},\mathcal{R}\mu ,\mathcal{I}\mu :III,1,5.\) \(\mathcal{M}(\mathrm{X};\mathbf{R}),\mathcal{M}(\mathrm{X}),\mathcal{M}_{+}(\mathrm{X}):III,1,5.\) \(\mu \leqslant \nu\) （ \(\mu ,\nu\) 为实测度）：III, 1, 5。 \(\mu^{+},\mu^{-},|\mu |\) （ \(\mu\) 为实测度）：III, 1, 5。 \(|\mu |\) （ \(\mu\) 为复测度）：III, 1, 6。 \(\| \mu \|\) （ \(\mu\) 为测度）：III, 1, 8。 \(\mathcal{M}^1 (\mathrm{X},\mathbf{R}),\mathcal{M}^1 (\mathrm{X}):III,1,8.\) \(\mu |Y\) （ \(\mu\) 为 X 上的测度，Y 为 X 的开子空间）：III, 2, 1。\\
Supp(μ)（ \(\mu\) 为测度）：III, 2, 2。 \(\langle \mathbf{f},\mathbf{z}'\rangle :III,3,1.\) \(\widetilde{\mathcal{K}} (\mathrm{X};\mathrm{E}):III,3,1.\) \(\int \mathbf{f}d\mu ,\int \mathbf{f}\mu ,\int \mathbf{f}(x)d\mu (x),\int \mathbf{f}(x)\mu (x)\) （ \(f\) 为 \(\widetilde{\mathcal{K}} (\mathrm{X};\mathrm{E}))\) 中的函数）：III, 3, 1。 \(\int d\mu (y)\int f(x,y)d\lambda (x):III,4,1.\) \(\iint f d\lambda d\mu ,\iint f d\mu d\lambda ,\iint f\lambda \mu ,\iint f\mu \lambda ,\iint f(x,y)d\lambda (x)d\mu (y),\) \(\iint f(x,y)d\mu (y)d\lambda (x),\iint f(x,y)\lambda (x)\mu (y),\iint f(x,y)\mu (y)\lambda (x):III,4,1.\)

\(\lambda \otimes \mu (\lambda ,\mu \text{为测度}):\mathrm{III},4,2.\) \(\mu_1\otimes \mu_2\otimes \dots \otimes \mu_n,\bigotimes_{i = 1}^n\mu_i:\mathrm{III},4,4.\) \(\int f d\mu_1d\mu_2\dots d\mu_n,\iint \dots \int f d\mu_1d\mu_2\dots d\mu_n,\int f(\mu_1\otimes \mu_2\otimes \dots \otimes \mu_n),\) \(\iint \dots \int f(x_1,x_2,\ldots ,x_n)d\mu_1(x_1)d\mu_2(x_2)\dots d\mu_n(x_n),\) \(\iint \dots \int f(x_1,x_2,\ldots ,x_n)\mu_1(x_1)\mu_2(x_2)\dots \mu_n(x_n):\mathrm{III},4,4.\) \(\bigotimes_{\lambda \in L}\mu_{\lambda}:\mathrm{III},4,6.\)

\(\varphi_{\mathrm{A}}:\mathrm{IV},1,1.\) \(\mathcal{K}_{+},\mathcal{I}_{+}(\mathrm{X}),\mathcal{I}_{+}:\mathrm{IV},1,1.\) \(\mu^{*}(f)\) （\(\mu\) 为正测度）：IV, 1, 1, IV, 1, 3 与 IV, 4, 习题 5。

\(\int^{*}fd\mu,\int^{*}f\mu,\int^{*}f(x)d\mu(x),\int^{*}f(x)\mu(x)\) （ \(f\) 为函数 \(\geqslant 0\) ， \(\mu\) 为正测度）：IV, 1, 3。

\(\mu^{*}(\mathrm{A})\) （A 为 X 的子集， \(\mu\) 为正测度）：IV, 1, 2 与 IV, 1, 4。

\(\widetilde{\mathbf{f}}:\mathrm{IV},2,4\) 与 IV, 2, 5.

\(\varphi((\widetilde{\mathbf{f}}_n))\), \(\widetilde{\mathbf{f}} +\widetilde{\mathbf{g}}\), \(\alpha \widetilde{\mathbf{f}},\widetilde{g}\widetilde{\mathbf{f}}:\mathrm{IV},2,4.\) \(\widetilde{f}\leqslant \widetilde{g},\sup_{n}\widetilde{f}_n,\inf_n\widetilde{f}_n,\limsup_{n\to \infty}\widetilde{f}_n,\liminf_{n\to \infty}\widetilde{f}_n\) （\(f,g,f_n\) 为数值函数）：IV, 2, 6。

\(|\mathbf{z}|\) （z 为赋范空间中的点）：IV, 3, 2。

\(|\mathbf{f}|\) （f 为取值于赋范空间的函数）：IV, 3, 2。

\(\mathrm{N}_p(\mathbf{f},\mu),\mathrm{N}_p(\mathbf{f}),\mathrm{N}_p(\widetilde{\mathbf{f}})\) （\(1\leqslant p<+\infty\)）：IV, 3, 2。

\(\mathcal{F}_{\mathrm{F}}(\mathrm{X}),\mathcal{F}_{\mathrm{F}}:\mathrm{IV},3,3.\) \(\mathcal{F}_{\mathrm{F}}^{p}(\mathrm{X},\mu),\mathcal{F}_{\mathrm{F}}^{p}(\mu),\mathcal{F}_{\mathrm{F}}^{p}:\mathrm{IV},3,3.\) \(\mathcal{N}_{\mathrm{F}}:\mathrm{IV},3,3.\) \(\mathcal{K}_{\mathrm{F}},\mathcal{L}_{\mathrm{F}}^{p}(\mathrm{X},\mu),\mathcal{L}_{\mathrm{F}}^{p}(\mu),\mathcal{L}_{\mathrm{F}}^{p},\mathrm{L}_{\mathrm{F}}^{p}(\mathrm{X},\mu),\mathrm{L}_{\mathrm{F}}^{p}(\mu),\mathrm{L}_{\mathrm{F}}^{p},\mathcal{L}^{p},\mathrm{L}^{p}\) （\(1\leqslant p<+\infty\)）：IV, 3, 4。

\(\| \widetilde{\mathbf{f}} \| _p\) （\(1\leqslant p < + \infty\)）：IV, 3, 4。

\(\mu (\mathbf{f}),\int \mathbf{f}d\mu ,\int \mathbf{f}(x)d\mu (x),\int \mathbf{f}\mu ,\int \mathbf{f}(x)\mu (x),\mu (\widetilde{\mathbf{f}})\) （f 为 \(\mu\) -可积的、取值于 Banach 空间的函数）：IV, 4, 1。

\(\mu (\mathrm{A})\) （A 为 \(\mu\) -可积集）：IV, 4, 5。

\(\mathcal{C}'(\mathrm{X};\mathbf{C}):\mathrm{IV},4,8.\) \(\mathcal{E}(\Phi),\mathcal{E}_{\mathrm{F}}(\Phi)\) （\(\Phi\) 为一个集环）：IV, 4, 9。

\(\mu_{*}(f)\) （\(f\) 为函数）：IV, 4, 习题 5。

\[
\mu_ {*} (\mathrm{A}) (\text { A a set }): \mathrm{IV}, 4, \text { Exer. } 7.
\]

\[
\int_ {\mathrm{A}} \textbf {f} d \mu , \int_ {\mathrm{A}} \textbf {f} \mu , \int_ {\mathrm{A}} ^ {*} f d \mu , \int_ {\mathrm{A}} ^ {*} f \mu : \mathrm{IV}, 5, 6.
\]

\[
\mathcal {S} (\mathrm{A}, \mu ; \mathrm{F}), \mathcal {S} _ {\mathrm{F}} (\mathrm{A}, \mu), \mathcal {S} _ {\mathrm{F}} (\mu), \mathcal {S} _ {\mathrm{F}}: \mathrm{IV}, 5, 11.
\]

\[
\boldsymbol {W} (\mathrm{V}, \mathrm{B}, \delta): \mathrm{IV}, 5, 11.
\]

\[
\mathrm{S} (\mathrm{A}, \mu ; \mathrm{F}), \mathrm{S} _ {\mathrm{F}} (\mathrm{A}, \mu), \mathrm{S} _ {\mathrm{F}} (\mu), \mathrm{S} _ {\mathrm{F}}: \mathrm{IV}, 5, 11.
\]

\[
\mathrm{M} _ {\infty} (f), m _ {\infty} (f): \mathrm{IV}, 6, 2.
\]

\[
\mathrm{N} _ {\infty} (\mathbf {f}): \mathrm{IV}, 6, 3.
\]

\[
\mathcal {L} _ {\mathrm{F}} ^ {\infty} (\mathrm{X}, \mu), \mathcal {L} _ {\mathrm{F}} ^ {\infty} (\mu), \mathcal {L} _ {\mathrm{F}} ^ {\infty}, \mathcal {N} _ {\mathrm{F}} ^ {\infty}: \mathrm{IV}, 6, 3.
\]

\[
\| \dot {\mathbf {f}} \| _ {\infty}, \mathrm{N} _ {\infty} (\dot {\mathbf {f}}): \mathrm{IV}, 6, 3.
\]

\[
\mathrm{L} _ {\mathrm{F}} ^ {\infty} (\mathrm{X}, \mu), \mathrm{L} _ {\mathrm{F}} ^ {\infty} (\mu), \mathrm{L} _ {\mathrm{F}} ^ {\infty}, \mathcal {L} ^ {\infty}, \mathrm{L} ^ {\infty}: \mathrm{IV}, 6, 3.
\]

\[
\mathbf {b} _ {\mu}: \mathrm{IV}, 7, 1.
\]

\[
\mathrm{Ch} _ {\mathcal {H}} (\mathrm{X}), \mathrm{Ch} (\mathrm{X}), \check {\mathrm{S}} _ {\mathcal {H}} (\mathrm{X}), \check {\mathrm{S}} (\mathrm{X}): \mathrm{IV}, 7, 3.
\]

\[
\mu^ {\bullet} (f), \mu^ {\bullet} (\mathrm{A}), \int^ {\bullet} f d \mu , \int^ {\bullet} f (t) d \mu (t), \int^ {\bullet} f \mu : \mathrm{V}, 1, 1.
\]

\[
\overline {{\mathcal {F}}} _ {\mathrm{F}} ^ {p} (\mathrm{T}, \mu), \overline {{\mathcal {F}}} _ {\mathrm{F}} ^ {p} (\mu), \overline {{\mathcal {F}}} _ {\mathrm{F}} ^ {p}: \mathrm{V}, 1, 3.
\]

\[
\overline {{\mathbf {N}}} _ {p} (f), \overline {{\mathcal {L}}} _ {\mathrm{F}} ^ {p} (\mathrm{T}, \mu), \overline {{\mathcal {L}}} _ {\mathrm{F}} ^ {p} (\mu), \overline {{\mathcal {L}}} _ {\mathrm{F}} ^ {p}: \mathrm{V}, 1, 3.
\]

\[
\overline {{\mathcal {L}}} _ {\mathrm{F}} ^ {p} (\mathrm{T}, \theta) (\theta \text {a complex measure}): \mathrm{V,1,3.}
\]

\[
\int \lambda_ {t} d \mu (t) (t \mapsto \lambda_ {t} \text {a family of positive measures}): \mathrm{V}, 3, 1.
\]

\[
\int d \mu (t) \int f (x) d \lambda_ {t} (x): \mathrm{V}, 3, 1.
\]

\[
\| \Lambda \| (\Lambda \text {a diffusion}): \mathrm{V}, 3, 5.
\]

\[
\langle \eta , h \rangle : \mathrm{V}, 3, 5.
\]

\[
\Lambda f, \mu \Lambda : \mathrm{V}, 3, 5.
\]

\[
\Lambda \mathrm{H}: \mathrm{V}, 3, 6.
\]

\[
\mathcal {L} _ {\mathrm{loc}} ^ {1} (\mathrm{T}, \mu ; \mathrm{F}), \mathrm{L} _ {\mathrm{loc}} ^ {1} (\mathrm{T}, \mu ; \mathrm{F}): \mathrm{V}, 5, 1.
\]

\[
\int_ {\mathrm{A}} ^ {\bullet} f d \mu : \mathrm{V}, 5, 3.
\]

\[
u (\mu_ {1}, \dots , \mu_ {n}) (u \text {a positively homogeneous numerical function}): \mathrm{V}, 5, 9.
\]

\[
\pi^ {- 1} (\mu) (\pi \text {a local homeomorphism}) \colon \mathrm{V,6,6.}
\]

\[
\iint^ {*} f (t, t ^ {\prime}) d \mu (t) d \mu^ {\prime} (t ^ {\prime}), \iint^ {\bullet} f (t, t ^ {\prime}) d \mu (t) d \mu^ {\prime} (t ^ {\prime}), \iint \mathbf {f} (t, t ^ {\prime}) d \mu (t) d \mu^ {\prime} (t ^ {\prime}): \mathrm{V}, 8, 1
\]

第 VI 章 :

\(\mathbf{F}^{\prime},\mathbf{F}^{\prime \prime},\mathbf{F}^{\prime *},\mathbf{F}_{\sigma}\) （F 为 Hausdorff 局部凸空间）：VI, 引言。

\(\mathcal{K}(\mathrm{T}),\mathcal{K}_{\mathbf{R}}(\mathrm{T}),\mathcal{K}_{\mathbf{C}}(\mathrm{T}),\mathcal{K}(\mathrm{T},\mathrm{A}),\mathcal{K}_{\mathbf{C}}(\mathrm{T},\mathrm{A}):\mathrm{VI},\) 引言。

\(\langle \mathbf{f},\mathbf{z}^{\prime}\rangle ,\langle \mathbf{z}^{\prime},\mathbf{f}\rangle :\mathrm{VI},1.\)

\(\int \mathbf{f} d\mu, \int \mathbf{f}(t) d\mu(t)\) （f 为向量值函数， \(\mu\) 为正测度）：VI, 1, 1。

\(g\mathbf{f},\mathbf{f}g\) （f 为向量值函数， \(g\) 为标量函数）：VI, 1, 1。

\(\mathcal{C}'(\mathrm{T}) : \mathrm{VI}, 1, 6.\)

\(\int f d\mathbf{m}, \int f(t) d\mathbf{m}(t)\) （ \(f\) 为数值函数， \(\mathbf{m}\) 为向量测度）：VI, 2, 1 与 VI, 2, 2。

\(g \cdot m\) （g 为数值函数，m 为向量测度）：VI, 2, 1。

\(\mathcal{L}(\mathbf{m}) : \mathrm{VI}, 2, 2.\)

\(q(\mathbf{m}), |\mathbf{m}|\) （ \(q\) 为半范数， \(\mathbf{m}\) 为向量测度）：VI, 2, 3。

\(\mathbf{f} \cdot \mu\) （f 为向量值函数， \(\mu\) 为正测度）：VI, 2, 4。

\(\mathcal{L}_{\mathrm{F}_s'}^{\infty}, \mathrm{L}_{\mathrm{F}_s'}^{\infty} : \mathrm{VI}, 2, 5.\)

\(\langle \mathbf{f},\mathbf{g}\rangle\) （f,g 为向量值函数）：VI, 2, 6。

\(\mathrm{I}_{\Phi ,\mathbf{m}},\int \mathbf{f}dm\) （f 为向量值函数，m 为向量测度）：VI, 2, 7。

\(|m|, \int \mathbf{f} dm \ (m \ a \ complex \ measure) : VI, 2, 8, III, 1, 6.\)

\(\mathcal{L}_{\mathrm{F}}^{p}(\mathrm{T},m), \overline{\mathcal{L}}_{\mathrm{F}}^{p}(\mathrm{T},m), \mathrm{L}_{\mathrm{F}}^{p}(\mathrm{T},m) (m \text{ 为复测度}): \mathrm{VI}, 2, 8, \mathrm{V}, 1, 3.\)

\(h \cdot m\) （ \(m\) 为复测度）：VI, 2, 8。

\(\overline{m}\) （ \(m\) 为复测度）：VI, 2, 8, III, 1, 5。

\(\| m\|\) （ \(m\) 为复测度）：VI, 2, 9, III, 1, 8。

\(\pi (m), m_{\mathrm{Y}}, m \otimes m' (m, m'\) 为复测度）：VI, 2, 10。

\(\mathfrak{B}(\mathrm{F}_{1},\mathrm{F}_{2}),{}^{r}\Phi,{}^{l}\Phi :\mathrm{VI},\mathrm{App.},1.\)

\(\mathrm{E}_{\sigma}, \mathrm{F}_{\sigma}, \mathrm{E}_{s}^{\prime}, \mathrm{F}_{s}^{\prime}, \mathcal{B}(\mathrm{E}, \mathrm{F}) : \mathrm{VI}, \mathrm{App.}, 1.\)

\(\Lambda_{\mathrm{F}^{\prime}}^{p}(\mathrm{T},\mu),\mathrm{M}_{p},\mathrm{M}_{p}^{\prime}:\mathrm{VI},1,\mathrm{Exer.}16.\)
% semantic-fragments: legacy-input-seam
\relax{}
